\documentclass[12pt, a4paper]{article}

% ---- 基础包 ----
\usepackage[UTF8]{ctex}
\usepackage[margin=2.5cm]{geometry}
\usepackage{listings}
\usepackage{xcolor}
\usepackage{hyperref}
\usepackage{titlesec}
\usepackage{parskip}
\usepackage{microtype}

% ---- 代码高亮配置 ----
\definecolor{codebg}{RGB}{245,245,245}
\definecolor{codeframe}{RGB}{180,180,180}
\definecolor{kwcolor}{RGB}{0,0,180}
\definecolor{strcolor}{RGB}{163,21,21}
\definecolor{cmtcolor}{RGB}{0,128,0}
\definecolor{numcolor}{RGB}{128,0,128}

\lstdefinestyle{python}{
  language=Python,
  backgroundcolor=\color{codebg},
  frame=single,
  rulecolor=\color{codeframe},
  basicstyle=\ttfamily\small,
  keywordstyle=\color{kwcolor}\bfseries,
  stringstyle=\color{strcolor},
  commentstyle=\color{cmtcolor}\itshape,
  numberstyle=\tiny\color{numcolor},
  numbers=left,
  numbersep=8pt,
  breaklines=true,
  breakatwhitespace=false,
  showstringspaces=false,
  tabsize=4,
  captionpos=b,
  xleftmargin=1.5em,
  framexleftmargin=1.5em,
  literate={->}{{\textrightarrow}}{2}
           {>=}{{\(\geq\)}}{2}
           {<=}{{\(\leq\)}}{2},
}

\lstdefinestyle{shell}{
  backgroundcolor=\color{codebg},
  frame=single,
  rulecolor=\color{codeframe},
  basicstyle=\ttfamily\small,
  breaklines=true,
  showstringspaces=false,
  xleftmargin=1.5em,
  framexleftmargin=1.5em,
}

\lstset{style=python}

% ---- 标题样式 ----
\titleformat{\section}{\large\bfseries}{}{0em}{}[\titlerule]
\titleformat{\subsection}{\normalsize\bfseries}{}{0em}{}

% ---- 文档信息 ----
\title{\textbf{密码学实验一：经典密码自动攻击与分析}\\[4pt]
       \large 关键实现步骤详解}
\date{}

% =============================================================
\begin{document}
\maketitle

本实验围绕三种经典密码——凯撒密码、单表替换密码、维吉尼亚密码——实现自动化的密钥搜索与明文还原。整体采用模块化设计，将通用工具集中在 \texttt{common.py} 中，三种攻击逻辑分别放入独立模块，最后由 \texttt{main.py} 统一调用和演示。这样做的好处是每个模块可以独立测试，公共评分逻辑只需维护一份，新增攻击方法时也不影响已有代码。

% =============================================================
\section{一、通用步骤（\texttt{common.py}）}

\texttt{common.py} 是整个项目的地基。所有攻击模块都从这里导入工具函数，因此它的设计质量直接决定整个实验的可维护性和评分准确性。

\subsection{1. 文本预处理}

在对密文进行任何分析之前，必须先统一文本的格式。英文密码学中通常约定一切字母均以大写处理，且只关心字母本身，标点和数字不参与加解密运算。如果密文中混有小写、数字或标点，直接参与计算会导致映射关系出错，因此预处理是不可省略的第一步。

代码中实现了两个预处理函数，共同完成这项工作：

\begin{lstlisting}[style=python, caption={文本预处理函数（common.py）}]
def normalize_text(text: str, keep_spaces: bool = True) -> str:
    text = text.upper()
    if keep_spaces:
        text = re.sub(r"[^A-Z\s]", " ", text)
        return re.sub(r"\s+", " ", text).strip()
    return re.sub(r"[^A-Z]", "", text)

def letters_only(text: str) -> str:
    return normalize_text(text, keep_spaces=False)
\end{lstlisting}

\texttt{normalize\_text} 先将全文转为大写，然后根据参数决定是否保留空格。保留空格时，会用正则把所有非字母非空格字符替换为空格，再压缩连续空格；不保留时，则直接去掉所有非字母字符，仅剩纯字母序列。\texttt{letters\_only} 是对后者的快捷封装。

设计两个版本的原因在于：评分时需要两种不同形态的文本。频率分析需要看连续字母序列的分布，必须去掉空格；而单词命中判断则需要将文本按空格分割为单词列表。两个函数分别覆盖这两种需求，调用时按需选择，不用在评分函数内部重复处理。

\subsection{2. 构造评分函数}

密码破解的自动化最终依赖于一个关键问题：如何判断一段解密结果``像不像英文''？这个问题没有绝对答案，因此代码采用了多维度加权的评分策略，将四种独立的统计信号融合为一个综合得分。

\textbf{第一个信号是字母频率。}英文中各字母出现的频率有固定的统计规律，E 约占 12.7\%，T 约占 9.1\%，而 Q 和 Z 则极为罕见。如果一段文本的字母频率分布与英文统计高度吻合，它是真实英文的概率就更高。代码使用卡方检验量化这种吻合程度：

\begin{lstlisting}[style=python, caption={英文字母频率表与频率评分函数（common.py）}]
ENGLISH_FREQ = {
    "A": 8.167, "B": 1.492, "C": 2.782, "D": 4.253, "E": 12.702,
    "F": 2.228, "G": 2.015, "H": 6.094, "I": 6.966, "J": 0.153,
    "K": 0.772, "L": 4.025, "M": 2.406, "N": 6.749, "O": 7.507,
    "P": 1.929, "Q": 0.095, "R": 5.987, "S": 6.327, "T": 9.056,
    "U": 2.758, "V": 0.978, "W": 2.360, "X": 0.150, "Y": 1.974,
    "Z": 0.074,
}

def frequency_score(text: str) -> float:
    cleaned = letters_only(text)
    if not cleaned:
        return float("-inf")
    total = len(cleaned)
    counts = Counter(cleaned)
    chi_square = 0.0
    for letter in ALPHABET:
        observed = counts.get(letter, 0)
        expected = total * ENGLISH_FREQ[letter] / 100
        if expected > 0:
            chi_square += ((observed - expected) ** 2) / expected
    return -chi_square
\end{lstlisting}

卡方值越小，说明实际分布与英文期望分布越接近，函数取负数返回，使得``更像英文''对应``得分越高''，统一了后续排序的方向。

\textbf{第二个信号是常用单词命中。}字母频率只反映单个字母的分布，无法判断字母组合是否构成有意义的英文单词。为此，代码维护了一个常用英文单词集合 \texttt{COMMON\_WORDS}，并实现 \texttt{word\_score} 对候选明文中每个单词逐一检查：

\begin{lstlisting}[style=python, caption={常用单词集合与单词评分函数（common.py）}]
COMMON_WORDS = {
    "THE", "BE", "TO", "OF", "AND", "A", "IN", "THAT", "HAVE", "I",
    "IT", "FOR", "NOT", "ON", "WITH", "HE", "AS", "YOU", "DO", "AT",
    "THIS", "BUT", "HIS", "BY", "FROM", "THEY", "WE", "SAY", "HER",
    "SHE", "OR", "AN", "WILL", "MY", "ONE", "ALL", "WOULD", "THERE",
    "THEIR", "IS", "ARE", "AM", "STUDENT", "LOVE", "QUICK", "BROWN",
    "FOX", "JUMPS", "OVER", "LAZY", "DOG", "NICE", "BOY", "HELLO",
    "WORLD", "ATTACK", "DAWN", "KEY", "SECRET", "MESSAGE", "TEST"
}

def word_score(text: str) -> float:
    words = normalize_text(text, keep_spaces=True).split()
    if not words:
        return 0.0
    score = 0.0
    for word in words:
        if word in COMMON_WORDS:
            score += 6.0 + len(word) * 0.3
        elif len(word) == 1 and word in {"A", "I"}:
            score += 2.0
        elif len(word) >= 3 and any(token in word for token in ("THE", "ING", "ION")):
            score += 1.5
    return score
\end{lstlisting}

命中完整单词得分最高，且更长的单词额外加分（因为长词的偶然命中概率更低，提供的信息量更大）。单字母 A 和 I 也单独处理，因为它们是合法的英文单词但可能被随机出现。

\textbf{第三个信号是 n-gram 统计。}相邻字母的组合频率也是判断语言的重要依据。TH、HE、IN 等双字母组合在英文中极为常见，随机序列中出现的概率很低。代码维护了常用双字母组合集 \texttt{COMMON\_BIGRAMS} 和三字母组合集 \texttt{COMMON\_TRIGRAMS}：

\begin{lstlisting}[style=python, caption={n-gram 集合与 n-gram 评分函数（common.py）}]
COMMON_BIGRAMS = {
    "TH", "HE", "IN", "ER", "AN", "RE", "ON", "AT", "EN", "ND",
    "TI", "ES", "OR", "TE", "OF", "ED", "IS", "IT", "AL", "AR",
    "ST", "TO", "NT", "NG", "SE", "HA", "AS", "OU", "IO", "LE"
}
COMMON_TRIGRAMS = {
    "THE", "AND", "ING", "HER", "ERE", "ENT", "THA", "NTH", "WAS",
    "ETH", "FOR", "HAT", "ION", "TIO", "VER"
}

def ngram_score(text: str) -> float:
    cleaned = letters_only(text)
    bigrams  = sum(1 for i in range(len(cleaned) - 1)
                   if cleaned[i:i+2] in COMMON_BIGRAMS)
    trigrams = sum(1 for i in range(len(cleaned) - 2)
                   if cleaned[i:i+3] in COMMON_TRIGRAMS)
    return bigrams * 1.2 + trigrams * 2.0
\end{lstlisting}

三字母组合权重高于双字母组合（2.0 vs 1.2），原因是三字母组合更稀有，出现则更有说服力。

\textbf{第四个信号是元音比例。}英文中元音（AEIOU）占字母总数的比例约为 38\%。纯辅音序列或全元音序列在正常英文中几乎不出现，因此偏离 38\% 越远则扣分越多：

\begin{lstlisting}[style=python, caption={元音比例评分函数（common.py）}]
def vowel_ratio_score(text: str) -> float:
    cleaned = letters_only(text)
    if not cleaned:
        return 0.0
    ratio = sum(1 for ch in cleaned if ch in "AEIOU") / len(cleaned)
    return -abs(ratio - 0.38) * 20
\end{lstlisting}

最终，四个子分以不同权重加权合并，得到综合评分函数：

\begin{lstlisting}[style=python, caption={综合评分函数（common.py）}]
def score_plaintext(text: str) -> float:
    return (
        frequency_score(text)     * 0.45
        + word_score(text)        * 1.8
        + ngram_score(text)       * 1.5
        + vowel_ratio_score(text) * 0.6
    )
\end{lstlisting}

权重的设定经过调试校准。单词命中权重最高（1.8），因为出现完整有意义的英文单词是最强的``这是真实明文''的信号；n-gram 次之（1.5），提供结构连贯性判据；频率得分和元音比例权重较低，作为辅助约束。这种多维度融合的设计比单一使用字母频率更鲁棒，能在密文较短时仍给出合理的排名。

\subsection{3. 记录候选结果}

每次解密尝试得到一个候选结果，代码以字典形式统一记录三项信息：密钥（凯撒密码中为整数移位量，维吉尼亚中为字符串密钥，单表替换中为字母映射字典）、解密后的明文字符串、以及调用 \texttt{score\_plaintext} 得到的浮点评分。

为了在实验报告中能够比较不同攻击方案的效率，还需要记录运行耗时。代码提供了 \texttt{timed\_call} 工具函数，内部使用 \texttt{time.perf\_counter} 实现微秒级精度的计时：

\begin{lstlisting}[style=python, caption={计时工具函数（common.py）}]
def timed_call(func: Callable, *args, **kwargs):
    start = time.perf_counter()
    result = func(*args, **kwargs)
    elapsed = time.perf_counter() - start
    return result, elapsed
\end{lstlisting}

调用时只需将目标函数和参数传入，即可同时获得函数返回值与运行时间，不需要在每个攻击模块里重复写计时逻辑。

\subsection{4. 排序输出}

所有候选收集完毕后，调用 \texttt{top\_results} 按得分降序排列并截取前若干条；格式化输出由 \texttt{format\_results} 完成，将结果列表转换为带序号的多行字符串，每行包含密钥、得分和明文片段，供 \texttt{main.py} 直接打印，方便人工逐条对比判断：

\begin{lstlisting}[style=python, caption={排序与格式化输出函数（common.py）}]
def top_results(results: Sequence[Dict], top_n: int = 5) -> List[Dict]:
    return sorted(results, key=lambda item: item["score"], reverse=True)[:top_n]

def format_results(results: Iterable[Dict]) -> str:
    lines = []
    for index, item in enumerate(results, start=1):
        lines.append(
            f"{index}. key={item['key']} | score={item['score']:.2f}"
            f" | text={item['plaintext']}"
        )
    return "\n".join(lines)
\end{lstlisting}

% =============================================================
\section{二、凯撒加密的关键步骤（\texttt{caesar\_attack.py}）}

凯撒密码是最简单的位移密码，整个加密过程只有一个参数：把明文中每个字母在字母表中向右循环位移若干位，解密时反向位移相同的量即可还原。密钥空间仅为 0 到 25 共 26 种，这是密码学里最小的密钥空间之一，穷举攻击几乎没有难度。

\subsection{1. 单字符移位规则与解密函数}

解密的核心思路是：已知密文字母在字母表中的位置，减去移位量后取模 26，得到明文字母的位置。取模运算保证结果始终落在 0 到 25 之间，从而实现字母表的循环环绕。非字母字符（空格、标点）不参与运算，原样保留，以维持密文的可读格式：

\begin{lstlisting}[style=python, caption={凯撒解密函数（caesar\_attack.py）}]
def decrypt_caesar(ciphertext: str, shift: int) -> str:
    output = []
    for ch in ciphertext.upper():
        if ch in ALPHABET:
            index = (ALPHABET.index(ch) - shift) % 26
            output.append(ALPHABET[index])
        else:
            output.append(ch)
    return "".join(output)
\end{lstlisting}

以具体例子说明：密文字符 D 在字母表中位于第 3 位（从 0 计数），若移位量 \texttt{shift=3}，则 $(3-3)\bmod 26=0$，对应字母表第 0 位即 A，正确还原。若 \texttt{shift} 大于字母位置，例如密文字符 B（位置 1）、\texttt{shift=3}，则 $(1-3)\bmod 26=24$，对应字母 Y，同样正确处理了负数情况，体现了取模运算在密码学中的重要作用。

\subsection{2. 遍历全部密钥并评分}

凯撒攻击的主函数逻辑非常简洁：循环 26 次，每次使用不同的移位量解密，调用评分函数，最后排序取最优：

\begin{lstlisting}[style=python, caption={凯撒穷举攻击函数（caesar\_attack.py）}]
def brute_force_caesar(ciphertext: str, top_n: int = 5) -> List[Dict]:
    results = []
    for shift in range(26):
        plaintext = decrypt_caesar(ciphertext, shift)
        results.append({
            "key": shift,
            "plaintext": plaintext,
            "score": score_plaintext(plaintext)
        })
    return top_results(results, top_n)
\end{lstlisting}

整体时间复杂度为 $O(26 \times N)$，其中 $N$ 为密文长度。即使密文很长，26 次线性扫描也可在毫秒内完成。由于密钥空间极小，这是密码学中少数可以完美穷举的情形。

\subsection{3. 保留高分结果与判断}

26 个候选结果按评分排序后，真实移位量对应的候选通常得分明显高于其他 25 个。这是因为正确解密后的文本符合英文字母频率、包含常见单词、n-gram 连贯，四个子分均会给出较高值，加权后综合分远超随机乱码对应的候选。在 \texttt{main.py} 的演示中，密文 ``WKHTXLFN EURZQ IRA MXPSV RYHU WKH ODCB GRJ'' 在 \texttt{shift=3} 时得分最高，还原为 ``THE QUICK BROWN FOX JUMPS OVER THE LAZY DOG''，与预期完全一致。人工审查时只需确认第一名候选的内容是否为有意义的英文，即可完成判断。

% =============================================================
\section{三、单表替换加密的关键步骤（\texttt{substitution\_attack.py}）}

单表替换密码将每个字母映射到另一个固定字母，整个字母表构成一个排列，密钥是这个排列本身。与凯撒不同，它的密钥空间是 $26! \approx 4.03\times10^{26}$，比凯撒大约 $10^{24}$ 倍，这从根本上决定了不能像凯撒那样直接穷举所有可能。

\subsection{1. 复杂度分析：为什么完整穷举不可行}

代码专门实现了一个函数来计算并展示这个结论：

\begin{lstlisting}[style=python, caption={26!复杂度计算函数（substitution\_attack.py）}]
def factorial_26_description() -> str:
    value = 1
    for number in range(2, 27):
        value *= number
    return f"26! = {value} ~= 4.03e26"
\end{lstlisting}

$26!$ 等于 403291461126605635584000000，约为 $4.03\times10^{26}$。假设计算机每秒能验证 $10^{12}$ 种映射（这已是极为乐观的估计），穷举完所有可能仍需约 $4\times10^{14}$ 年，远超宇宙年龄（约 $1.4\times10^{10}$ 年）。实验报告的结论部分必须明确写出：完整 26 字母单表替换密码的纯穷举不可行，需要借助约束或启发式方法。

正因如此，代码提供了两种互补的攻击方案：一是缩小字母集的局部穷举，二是利用词典约束的束搜索。

\subsection{2. 方案一：缩小字母集穷举}

如果能事先确定密文中只出现了少量不同字母（例如已知明文攻击场景、或短语加密），则可以将穷举范围限制在这些字母上，将 $26!$ 的复杂度降低到 $k!$（$k$ 为实际出现的字母数）。$7!=5040$，$10!=3628800$，这样的规模是完全可以在秒级内穷举的。

函数接收密文、密文字母集、明文候选字母集三个参数，用 \texttt{itertools.permutations} 生成明文字母集的所有排列，每个排列与密文字母集按位组合成映射字典，再将映射作用于密文，调用评分函数筛选：

\begin{lstlisting}[style=python, caption={缩小字母集穷举函数（substitution\_attack.py）}]
def apply_substitution_mapping(ciphertext: str, mapping: Dict[str, str]) -> str:
    output = []
    for ch in ciphertext.upper():
        if "A" <= ch <= "Z":
            output.append(mapping.get(ch, "_"))
        else:
            output.append(ch)
    return "".join(output)

def brute_force_reduced_alphabet(
    ciphertext: str,
    cipher_letters: Sequence[str],
    plain_letters: Sequence[str],
    top_n: int = 10
) -> List[Dict]:
    if len(cipher_letters) != len(plain_letters):
        raise ValueError("cipher_letters and plain_letters must have the same length")
    results = []
    for perm in itertools.permutations(plain_letters):
        mapping = dict(zip(cipher_letters, perm))
        plaintext = apply_substitution_mapping(ciphertext, mapping)
        results.append({
            "key": mapping,
            "plaintext": plaintext,
            "score": score_plaintext(plaintext)
        })
    return top_results(results, top_n)
\end{lstlisting}

\texttt{apply\_substitution\_mapping} 将映射应用于密文时，对每个字母查找映射字典，若找不到则填充占位符 \texttt{"\_"}，避免因映射不完整导致错误结果混入。在 \texttt{main.py} 的演示中，密文 ``I LOVE YOU'' 仅涉及 I、L、O、V、E、Y、U 七个字母，穷举 5040 种排列可以秒级完成，并正确命中原始映射。

\subsection{3. 方案二：带词典约束的束搜索（回溯搜索）}

当密文是完整的英文句子且保留了单词边界（有空格）时，可以利用英文词典对搜索空间进行大幅剪枝。核心思路是：逐个处理密文中的每个单词，只尝试与该密文单词结构（长度和字母重复模式）相匹配的词典单词作为候选，并在扩展映射时立即检查一致性，不一致则直接丢弃该分支。

映射的表示方式是双向字典，同时维护正向映射（密文字母 $\to$ 明文字母）和反向映射（明文字母 $\to$ 密文字母）。双向结构使得一致性检查可以在 $O(1)$ 时间内完成，而不需要每次都遍历整个映射：

\begin{lstlisting}[style=python, caption={一致性检查的映射扩展函数（substitution\_attack.py）}]
def extend_mapping(current_map, reverse_map, cipher_word, plain_word):
    if len(cipher_word) != len(plain_word):
        return None
    next_map = dict(current_map)
    next_reverse = dict(reverse_map)
    for cipher_char, plain_char in zip(cipher_word, plain_word):
        existing_plain  = next_map.get(cipher_char)
        existing_cipher = next_reverse.get(plain_char)
        # 同一密文字母已映射到别的明文字母，冲突
        if existing_plain is not None and existing_plain != plain_char:
            return None
        # 同一明文字母已被其他密文字母占用，冲突
        if existing_cipher is not None and existing_cipher != cipher_char:
            return None
        next_map[cipher_char] = plain_char
        next_reverse[plain_char] = cipher_char
    return next_map, next_reverse
\end{lstlisting}

若发现任何冲突，立即返回 \texttt{None}，该分支被剪掉，不再继续扩展。这保证了最终映射始终是合法的单射（没有重复映射，没有冲突映射）。

为了进一步减少无效尝试，代码引入了单词模式分桶机制。\texttt{word\_pattern} 将单词转换为数字序列，例如 HELLO $\to$ (0,1,2,2,3)，DEED $\to$ (0,1,1,0)。具有相同重复结构的单词模式相同。算法预先将词典按（长度，模式）分桶，搜索时只需从匹配的桶中取候选，大幅过滤掉不可能成立的单词：

\begin{lstlisting}[style=python, caption={单词模式与束搜索主函数（substitution\_attack.py）}]
def word_pattern(word: str) -> tuple:
    indices = {}
    pattern = []
    next_index = 0
    for ch in word:
        if ch not in indices:
            indices[ch] = next_index
            next_index += 1
        pattern.append(indices[ch])
    return tuple(pattern)

def backtracking_with_dictionary(
    ciphertext: str,
    candidate_words=None,
    top_n: int = 10,
    beam_width=None,
) -> List[Dict]:
    words = [word for word in ciphertext.upper().split() if word]
    source_words = candidate_words if candidate_words else COMMON_WORDS
    dictionary = sorted({w.upper() for w in source_words if w.isalpha()})
    if beam_width is None:
        beam_width = max(top_n * 10, 100)

    # 按（长度, 模式）分桶
    candidate_buckets = {}
    for word in dictionary:
        bucket_key = (len(word), word_pattern(word))
        candidate_buckets.setdefault(bucket_key, []).append(word)

    beam = [{"key": {}, "reverse_key": {}, "plain_words": [], "score": 0.0}]
    for cipher_word in words:
        bucket_key = (len(cipher_word), word_pattern(cipher_word))
        candidates = candidate_buckets.get(bucket_key, [])
        if not candidates:
            return []
        expanded = {}
        for state in beam:
            for plain_word in candidates:
                extended = extend_mapping(
                    state["key"], state["reverse_key"],
                    cipher_word, plain_word
                )
                if extended is None:
                    continue
                next_map, next_reverse = extended
                next_plain_words = state["plain_words"] + [plain_word]
                score = frequency_score(" ".join(next_plain_words))
                signature = tuple(sorted(next_map.items()))
                existing = expanded.get(signature)
                if existing is None or score > existing["score"]:
                    expanded[signature] = {
                        "key": next_map, "reverse_key": next_reverse,
                        "plain_words": next_plain_words, "score": score,
                    }
        if not expanded:
            return []
        beam = sorted(expanded.values(),
                      key=lambda x: x["score"], reverse=True)[:beam_width]

    results = []
    for state in beam:
        plaintext = apply_substitution_mapping(ciphertext, state["key"])
        if "_" in plaintext:
            continue
        results.append({
            "key": dict(sorted(state["key"].items())),
            "plaintext": plaintext,
            "score": frequency_score(plaintext),
        })
    return top_results(results, top_n)
\end{lstlisting}


% =============================================================
\section{四、维吉尼亚加密的关键步骤（\texttt{vigenere\_attack.py}）}

维吉尼亚密码是凯撒密码的多表推广：使用一个关键字，对明文的每个字母依次用关键字中对应位置的字母作为移位量进行凯撒加密，关键字循环使用。例如关键字为 KEY（K=10，E=4，Y=24），则明文第 1、4、7\ldots 个字母用 10 位移加密，第 2、5、8\ldots 个用 4 位移，第 3、6、9\ldots 个用 24 位移。由于同一明文字母在不同位置会加密为不同密文字母，单纯的字母频率分析失效，破解难度远高于凯撒密码。

密钥空间为 $26^L$（$L$ 为密钥长度）：长度为 1 时与凯撒等价（26 种），长度为 2 时有 676 种，长度为 3 时有 17576 种，长度为 4 时超过 45 万种，之后每增加一位扩大 26 倍。这种指数增长决定了必须对密钥长度设置上限，否则搜索空间将迅速变得不可控。

\subsection{1. 设置最大密钥长度}

代码在 \texttt{brute\_force\_vigenere} 函数中明确设置了上限参数 \texttt{max\_key\_len}，并在参数大于 3 时直接抛出异常，阻止实验者误设过大的值：

\begin{lstlisting}[style=python, caption={维吉尼亚穷举函数入口（vigenere\_attack.py）}]
def brute_force_vigenere(
    ciphertext: str, max_key_len: int = 3, top_n: int = 10
) -> List[Dict]:
    if max_key_len < 1:
        raise ValueError("max_key_len must be at least 1")
    if max_key_len > 3:
        raise ValueError(
            "max_key_len above 3 is intentionally blocked "
            "in this lab demo to avoid explosive brute-force cost"
        )
    ...
\end{lstlisting}

这并非技术上的限制，而是一种实验安全保护。长度为 4 时需要穷举 456976 个密钥，乘以密文长度后运算量显著增大，不适合在课堂演示环境中不受限地运行。实验报告需记录 \texttt{max\_key\_len} 为 1、2、3 时的耗时对比，直观感受指数增长的影响。

\subsection{2. 实现加解密函数}

加密函数 \texttt{encrypt\_vigenere} 和解密函数 \texttt{decrypt\_vigenere} 互为对称操作，结构完全相同，区别仅在于位移方向：

\begin{lstlisting}[style=python, caption={维吉尼亚加解密函数（vigenere\_attack.py）}]
def encrypt_vigenere(plaintext: str, key: str) -> str:
    key = "".join(ch for ch in key.upper() if ch in ALPHABET)
    if not key:
        raise ValueError("key must contain at least one alphabetic character")
    output = []
    key_index = 0
    for ch in plaintext.upper():
        if ch in ALPHABET:
            shift = ALPHABET.index(key[key_index % len(key)])
            cipher_index = (ALPHABET.index(ch) + shift) % 26
            output.append(ALPHABET[cipher_index])
            key_index += 1
        else:
            output.append(ch)
    return "".join(output)

def decrypt_vigenere(ciphertext: str, key: str) -> str:
    key = "".join(ch for ch in key.upper() if ch in ALPHABET)
    if not key:
        raise ValueError("key must contain at least one alphabetic character")
    output = []
    key_index = 0
    for ch in ciphertext.upper():
        if ch in ALPHABET:
            shift = ALPHABET.index(key[key_index % len(key)])
            plain_index = (ALPHABET.index(ch) - shift) % 26
            output.append(ALPHABET[plain_index])
            key_index += 1
        else:
            output.append(ch)
    return "".join(output)
\end{lstlisting}

实现中有一个重要细节：\texttt{key\_index} 只在遇到字母时才递增，碰到空格或标点时跳过但不消耗关键字的轮次。这保证了加解密对非字母字符的处理是一致的，不会因标点符号导致关键字对齐位置偏移。

\texttt{main.py} 中先用 \texttt{encrypt\_vigenere("I AM A STUDENT", "KEY")} 加密，再用破解函数还原，以端到端验证加解密逻辑的正确性，避免因实现缺陷导致攻击函数评估的候选始终错误。

\subsection{3. 按长度从短到长穷举候选密钥}

穷举的核心使用 \texttt{itertools.product} 生成所有组合，外层循环密钥长度从 1 逐步到 \texttt{max\_key\_len}：

\begin{lstlisting}[style=python, caption={按长度穷举全部候选密钥（vigenere\_attack.py）}]
    results = []
    for key_len in range(1, max_key_len + 1):
        for key_tuple in itertools.product(ALPHABET, repeat=key_len):
            key = "".join(key_tuple)
            plaintext = decrypt_vigenere(ciphertext, key)
            results.append({
                "key": key,
                "plaintext": plaintext,
                "score": score_plaintext(plaintext)
            })
    return top_results(results, top_n)
\end{lstlisting}

\texttt{itertools.product(ALPHABET, repeat=key\_len)} 会生成 $26^{\texttt{key\_len}}$ 个元组，每个元组是长度为 \texttt{key\_len} 的字母序列，自动完成所有组合的枚举，无需手写嵌套循环。从短到长枚举的好处是：即使中途中断，已处理的较短密钥的结果仍然有效，可以查看部分结果。

\subsection{4. 评分与排序，以及时间复杂度对比}

每个密钥的候选明文都调用 \texttt{score\_plaintext} 评分，最终 \texttt{top\_results} 取前 \texttt{top\_n} 高分输出。正确密钥对应的解密结果应当在所有候选中得分最高，因为它是唯一一个能产生有意义英文的结果。

为了定量展示时间随密钥长度增长的变化，\texttt{main.py} 对每个攻击段落使用 \texttt{perf\_counter} 精确计时。以密文``I AM A STUDENT''用密钥 KEY 加密后的密文为例，实验中应分别记录 \texttt{max\_key\_len=1}、2、3 的耗时。设每次解密计算的时间为 $t$，则三种情况下总解密次数分别为 26、$26+676=702$、$26+676+17576=18278$ 次，耗时比约为 $1:27:703$，实验结果将直观地验证指数复杂度 $O(26^L \times N)$ 的规律。

% =============================================================
\section{五、主控模块（\texttt{main.py}）}

\texttt{main.py} 是整个实验的入口，负责将三种攻击整合为一个统一的可运行程序，支持演示模式和命令行模式两种调用路径。其设计思路是：不重复任何攻击逻辑，只负责参数解析、流程调度、计时记录和格式化输出，真正做到了主控层与业务层的分离。

\subsection{1. 演示模式（\texttt{verify\_cases}）}

不传任何命令行参数时，程序执行 \texttt{verify\_cases} 函数，依次运行若干预置测试用例。凯撒演示部分测试两组密文，单表替换演示部分先打印 $26!$ 量级再分别演示两种方案，维吉尼亚演示部分先加密再破解以端到端验证正确性。所有用例的计时都在 \texttt{main.py} 层面通过 \texttt{perf\_counter} 完成，最后汇总总运行时间。

\subsection{2. 命令行模式（\texttt{build\_parser} 与 \texttt{run\_cli}）}

命令行模式允许用户传入自定义密文并选择攻击方案，通过 \texttt{argparse} 实现参数解析：

\begin{lstlisting}[style=python, caption={命令行参数定义（main.py）}]
parser.add_argument("--ciphertext",  help="待破解密文")
parser.add_argument("--top-k",       type=int, default=5)
parser.add_argument("--run-caesar",  action="store_true")
parser.add_argument("--run-vigenere", action="store_true")
parser.add_argument("--max-key-len", type=int, default=3)
parser.add_argument("--run-substitution-reduced",    action="store_true")
parser.add_argument("--cipher-letters")
parser.add_argument("--plain-letters")
parser.add_argument("--run-substitution-dictionary", action="store_true")
parser.add_argument("--words",       nargs="+")
parser.add_argument("--beam-width",  type=int)
\end{lstlisting}

调用示例如下，可以灵活组合使用：

\begin{lstlisting}[style=shell, caption={命令行调用示例}]
# 对自定义密文运行凯撒穷举，输出前 3 个候选
python main.py --ciphertext "KHOOR ZRUOG" --run-caesar --top-k 3

# 对自定义密文运行词典束搜索
python main.py --ciphertext "A NICE BOY" --run-substitution-dictionary

# 对维吉尼亚密文穷举，密钥长度上限为 2
python main.py --ciphertext "MIAU" --run-vigenere --max-key-len 2
\end{lstlisting}

\subsection{3. 参数合法性验证（\texttt{validate\_args}）}

在执行攻击前，\texttt{validate\_args} 函数会检查参数的合法性，发现问题时直接调用 \texttt{parser.error} 打印清晰的错误提示并退出，避免程序运行到中途才因非法状态崩溃。检查项包括：必须至少选择一种攻击模式；\texttt{--ciphertext} 不得为空；\texttt{--top-k} 和 \texttt{--beam-width} 必须为正整数；缩小字母集模式要求 \texttt{--cipher-letters} 和 \texttt{--plain-letters} 等长且各自不含重复字母（保证映射为合法单射）。\texttt{main} 函数的整体调度逻辑非常清晰：解析参数，若没有任何攻击模式被指定则进入演示模式，否则进行合法性验证后进入命令行模式。这种``无参数走演示，有参数走 CLI''的设计使得程序既适合课堂演示，也适合对特定密文进行独立分析。

% =============================================================
\section{六、模块依赖关系与运行方法}

\texttt{common.py} 是基础模块，不依赖任何项目内其他文件，仅使用 Python 标准库（\texttt{re}、\texttt{time}、\texttt{collections}、\texttt{typing}）。三个攻击模块均从 \texttt{common.py} 导入所需的常量和函数：\texttt{caesar\_attack.py} 导入 \texttt{ALPHABET}、\texttt{score\_plaintext}、\texttt{top\_results}；\texttt{vigenere\_attack.py} 同样导入这三项；\texttt{substitution\_attack.py} 还额外导入了 \texttt{COMMON\_WORDS} 和 \texttt{frequency\_score}（束搜索中间状态使用频率得分，最终结果使用综合得分）。\texttt{main.py} 处于最顶层，从所有四个模块导入对外接口函数，自身不实现任何业务逻辑。

这种分层结构的优势在于：修改评分策略只需改动 \texttt{common.py}，所有攻击模块立即受益；新增一种密码攻击只需新建一个模块并在 \texttt{main.py} 中添加调用，不影响现有代码；单元测试时可以直接导入单个攻击模块进行验证，无需运行整个程序。

运行时，进入 \texttt{crypto\_lab/实验一/} 目录后直接执行 \texttt{python main.py} 即可触发演示模式，观察三种攻击在预置测试用例上的完整运行过程与计时结果。如需对特定密文进行分析，则按上文所示命令行格式附加相应参数即可。

\end{document}
